\chapter{Introduction}
\label{ch:intro}

\section{Ranking forecasters with dependent errors}
\label{sec:intro-motivation}

This thesis constructs simultaneous confidence sets for the ranks of forecasters whose
losses are correlated over time and across forecasters. Given estimated mean losses and
an estimate of their covariance, the procedure of \citet{mogstad2024} returns a
confidence set for every forecaster's rank, such that all true ranks lie in their sets
with probability at least $1-\alpha$ in large samples. I use that procedure unchanged. My contribution
is the covariance it needs: an estimate that is robust to serial dependence and that
does not reintroduce the common shocks all forecasters share.

How that covariance is built changes the answer. I develop two constructions. The first
reconstructs a covariance of the levels from the pairwise standard errors; the second
estimates the covariance of the pairwise contrasts directly. On a balanced panel they
give the same test, exactly under a common bandwidth and asymptotically otherwise. On the
unbalanced panels typical of forecast surveys they do not.
For the five best-covered nominal-output forecasters of the U.S.\ Survey of Professional
Forecasters, the reconstruction returns the uninformative set $[1,5]$ for every
forecaster, while the direct estimator separates three of them (Table~\ref{tab:ngdp}).

Rankings of forecasters are published routinely: central banks, survey providers and
the financial press order professional forecasters by average accuracy.%
\footnote{For example, the Banco Central do Brasil publishes monthly and annual ``Top~5''
rankings of the institutions in its Focus survey by forecast accuracy \citep{bcbtop5};
the Lawrence R.~Klein Award goes to the most accurate panelist in the \emph{Blue Chip
Economic Indicators} survey over the preceding four years \citep{bluechipklein2025}; and
Bloomberg periodically ranks economists by the accuracy of their forecasts of monthly
macroeconomic indicators \citep{bloomberg2017}.}
Such a ranking is a statistic. When the leading forecasters' mean squared errors differ
by less than the sampling noise, a different sample could reorder them, and the
forecaster ranked first may be indistinguishable from the one ranked fifth. A confidence
set for each rank states which orderings the data rule out.

\section{Two kinds of dependence}
\label{sec:intro-dependence}

Forecast losses are dependent in two ways. They are serially dependent: slow-moving
regimes, persistent models and overlapping information sets make a forecaster's errors
persist, so $n$ periods carry less information than $n$ independent observations.
Ignoring this understates the uncertainty of each mean loss and separates forecasters
who should not be separated. They are also cross-sectionally dependent: a recession or a
pandemic raises every forecaster's loss in the same quarter. Treating forecasters as
independent adds this common variance to every contrast and hides real differences in
skill.

I handle the two in different places. Serial dependence is a property of each pairwise
difference series, so I absorb it into a long-run (HAC) standard error of each mean
difference. Cross-sectional dependence is handled by working with differences of losses,
in which an additive common shock cancels, and by building the simulation covariance from
these differences rather than from levels. The difficulty is that the procedure's
critical value needs a full covariance matrix, not only the pairwise standard errors,
and that matrix must agree with the pairwise standard errors without bringing the
common shock back.

\section{Contribution}
\label{sec:intro-contribution}

The contribution has five parts.

\begin{enumerate}[nosep]
  \item \emph{Serial dependence} (Chapter~\ref{ch:serial}). I compute a Newey--West
  standard error for each difference series, with the Bartlett kernel and the
  \citet{andrews1991} plug-in bandwidth. The bandwidth adapts to each pair's
  persistence; a fixed-rate rule truncates too early on strongly persistent series,
  which makes the sets too narrow.
  \item \emph{Two covariance constructions} (Chapter~\ref{ch:cross}). The first
  reconstructs a level covariance from the pairwise standard errors by classical
  multidimensional scaling. The second estimates the covariance of the contrasts
  directly, with a pairwise-of-pairs version for unbalanced panels. I prove that the two
  coincide on a balanced panel --- exactly under a common bandwidth, asymptotically
  otherwise. On an unbalanced panel they estimate different objects, and only the direct
  estimator targets the covariance of the contrasts the procedure tests.
  \item \emph{Calibration} (\S\S\ref{sec:mds-psd}--\ref{sec:caveat} and
  \S\ref{sec:mc-mechanism}). The direct construction studentizes the simulated contrasts
  by its own diagonal, so they have unit variance. The reconstruction studentizes them by
  the original standard errors; when these are far from Euclidean, its positive
  semi-definite repair inflates the simulated variances and hence the critical value.
  \item \emph{Monte Carlo evidence} (Chapter~\ref{ch:mc}). Both constructions attain
  nominal joint coverage on every design. They are close to equivalent on balanced
  stationary panels, and the direct construction is sharper only when imbalance and
  non-stationarity are combined. A joint block bootstrap serves as a covariance-free
  benchmark.
  \item \emph{Applications} (Chapter~\ref{ch:apps}). The procedure separates little among
  macroeconomic forecasters, finds a top tier among tournament forecasters, and separates
  competent forecasting methods from systematic failures in the M3 competition without
  naming a winner. In each case the width of the sets measures how much the data can
  rank.
\end{enumerate}

\section{Roadmap}
\label{sec:intro-roadmap}

Chapter~\ref{ch:background} reviews the related literature and Chapter~\ref{ch:framework}
sets up the notation. Chapters~\ref{ch:serial}--\ref{ch:downstream} develop the method:
the standard errors, the two covariance constructions, and the step from covariance to
rank confidence sets. Chapter~\ref{ch:mc} reports the Monte Carlo evidence,
Chapter~\ref{ch:apps} the applications, and Chapter~\ref{ch:conclusion} concludes.
